\documentclass[11pt,a4paper]{article}
\usepackage[italian]{babel}
\usepackage[margin=2.5cm]{geometry}
\usepackage{parskip}
\usepackage{amsmath, amssymb, amsthm}
\usepackage{booktabs}
\usepackage{array}
\usepackage{xcolor}
\usepackage{needspace}
\usepackage{enumitem}
\usepackage{listings}
\usepackage{tikz}
\usepackage[most]{tcolorbox}
\usepackage[hidelinks]{hyperref}
\usetikzlibrary{decorations.pathreplacing, shapes.geometric, arrows.meta, positioning}

% Niente vedove/orfane: i paragrafi non lasciano righe isolate a fine/inizio pagina
\widowpenalty=10000
\clubpenalty=10000

% Elenchi più compatti
\setlist{itemsep=2pt, parsep=0pt, topsep=4pt, partopsep=0pt}

% --- Riquadri con bordo nero, senza numero (si spezzano tra due pagine) ---
% Uso: \begin{definizione} ... \end{definizione}
%      \begin{teorema}[Nome] ... \end{teorema}  (il nome è facoltativo)
\tcbset{nota/.style={enhanced jigsaw, breakable, colback=white,
  colframe=black, boxrule=0.5pt, sharp corners}}
\NewTColorBox{definizione}{}{nota,
  before upper={\textbf{Definizione.}\ }}
\NewTColorBox{teorema}{o}{nota,
  before upper={\textbf{Teorema\IfValueT{#1}{ (#1)}.}\ }}
\NewTColorBox{proposizione}{o}{nota,
  before upper={\textbf{Proposizione\IfValueT{#1}{ (#1)}.}\ }}
\NewTColorBox{proprieta}{o}{nota,
  before upper={\textbf{Proprietà\IfValueT{#1}{ (#1)}.}\ }}
\NewTColorBox{assioma}{o}{nota, boxrule=1pt,
  before upper={\textbf{Assioma\IfValueT{#1}{ (#1)}.}\ }}

% --- Ambienti semplici (senza riquadro) ---
% Il titolo sta in cima al blocco, anche se il contenuto inizia con un riquadro o
% con codice; e non resta mai da solo in fondo a una pagina.
% Uso: \begin{esempio}[Titolo facoltativo] ... \end{esempio}
\newcommand{\newrunin}[2]{%
  \NewDocumentEnvironment{#1}{o}{%
    \par\Needspace{4\baselineskip}\addvspace{\medskipamount}\noindent
    \textbf{#2}\IfValueT{##1}{ (##1)}\textbf{.}\ \ignorespaces}%
  {\par\addvspace{\medskipamount}}}
\newrunin{esempio}{Esempio}
\newrunin{esercizio}{Esercizio}
\newrunin{osservazione}{Osservazione}

% V e F nelle tavole di verità
\newcommand{\V}{\textsf{V}}
\newcommand{\F}{\textsf{F}}

% --- Codice C ---
% Uso: \begin{codice} ... \end{codice}      codice C numerato
%      \begin{comando} ... \end{comando}    comandi da terminale
%      \begin{risultato} ... \end{risultato} output di un programma
%      \lstinline|printf("ciao");|           codice dentro al testo
\lstdefinestyle{c}{
  language=C,
  basicstyle=\ttfamily\small,
  keywordstyle=\color{blue}\bfseries,
  commentstyle=\color{gray}\itshape,
  stringstyle=\color{red!70!black},
  numbers=left,
  numberstyle=\tiny\color{gray},
  numbersep=8pt,
  columns=flexible,
  keepspaces=true,
  breaklines=true,
  showstringspaces=false,
  tabsize=4,
  morekeywords={include, define},
  % lettere accentate nei commenti
  literate={à}{{\`a}}1 {è}{{\`e}}1 {é}{{\'e}}1 {ì}{{\`i}}1
           {ò}{{\`o}}1 {ù}{{\`u}}1 {È}{{\`E}}1 {À}{{\`A}}1
}
\lstset{style=c}
\newtcblisting{codice}{listing only, nota, left=6mm,
  listing options={style=c}}
\newtcblisting{comando}{listing only, nota,
  listing options={basicstyle=\ttfamily\small, numbers=none, language={},
    columns=flexible, keepspaces=true, breaklines=true}}
\newtcblisting{risultato}{listing only, enhanced jigsaw, breakable,
  colback=black!5, colframe=black, boxrule=0.5pt, sharp corners,
  listing options={basicstyle=\ttfamily\small, numbers=none, language={},
    columns=flexible, keepspaces=true, breaklines=true}}

% --- Diagrammi di flusso (simboli standard) ---
\tikzset{
  terminale/.style = {ellipse, draw, minimum width=2.5cm, minimum height=0.9cm},
  io/.style        = {trapezium, trapezium left angle=70, trapezium right angle=110,
                      draw, minimum width=2.5cm, minimum height=0.9cm},
  processo/.style  = {rectangle, draw, minimum width=2.5cm, minimum height=0.9cm},
  decisione/.style = {diamond, draw, aspect=2, minimum width=2.5cm},
  freccia/.style   = {thick, -{Stealth}}
}

\title{Continuità e limite in un punto}
\author{Marco Cerbella}
\date{Lezione 9 -- 22 ottobre 2026}

\begin{document}
\maketitle

\section{Funzioni continue e limite per $x \to x_0$}

\begin{definizione}[Funzione continua in un punto]
Si dice che $f : D \to \mathbb{R}$ è \emph{continua} in $x_0 \in D$ se
\[
\forall \varepsilon > 0,\ \exists \delta > 0 \text{ tale che } \forall x \in D : |x - x_0| < \delta \text{ si ha } |f(x) - f(x_0)| < \varepsilon .
\]
\end{definizione}

La condizione $|f(x) - f(x_0)| < \varepsilon$ significa $f(x) \in I_\varepsilon(f(x_0))$, cioè $f(x_0) - \varepsilon < f(x) < f(x_0) + \varepsilon$.

\begin{definizione}[Limite per $x \to x_0$]
Sia $f : D \to \mathbb{R}$ definita in un intorno di $x_0 \in \mathbb{R}$ tranne eventualmente nel punto $x_0$. Si dice che $f$ \emph{ha limite} $\ell \in \mathbb{R}$ (o tende a $\ell$) per $x$ tendente a $x_0$ e si scrive $\displaystyle\lim_{x \to x_0} f(x) = \ell$ se
\[
\forall \varepsilon > 0,\ \exists \delta > 0 \text{ tale che } \forall x \in D \setminus \{x_0\} : |x - x_0| < \delta \text{ si ha } |f(x) - \ell| < \varepsilon .
\]
\end{definizione}

\begin{osservazione}
Il limite si calcola per i valori $x_0$ che sono \emph{punti di accumulazione}. Il punto $x_0$ può non appartenere al dominio, ma deve essere di accumulazione. Nella condizione del limite il punto $x_0$ \emph{non c'è} (si considera $x \in D \setminus \{x_0\}$).
\end{osservazione}

\section{Punti di accumulazione}

\begin{definizione}[Punto di accumulazione]
Dato un insieme $A$, si dice che $x_0$ è un \emph{punto di accumulazione} per $A$ se per ogni $\varepsilon > 0$ esiste un intorno di ampiezza $\varepsilon$ di $x_0$ tale che $(I_\varepsilon(x_0) \setminus \{x_0\}) \cap A \neq \emptyset$.
\end{definizione}

\begin{proposizione}
Un numero $x_0 \in \mathbb{R}$ è di accumulazione per un insieme $A \subseteq \mathbb{R}$ se e solo se in ogni intorno di $x_0$ cadono infiniti punti di $A$ distinti da $x_0$.
\end{proposizione}

\begin{esempio}
\begin{itemize}
    \item Se $A = [a, b]$ ogni punto $x_0 \in A$ è di accumulazione per $A$.
    \item Se $A = \,]a, b[$ tutti i punti di $A$ sono punti di accumulazione, ma lo sono anche il punto $a$ e il punto $b$.
    \item Se $a \in \mathbb{R}$ e $A = \{a\}$ allora $A$ non ha punti di accumulazione.
    \item Se $A = \left\{ x \in \mathbb{R} : \exists n \in \mathbb{N} : x = \tfrac{1}{n} \right\} = \left\{1, \tfrac{1}{2}, \tfrac{1}{3}, \tfrac{1}{4}, \dots \right\}$ allora $x_0 = 0$ è il solo punto di accumulazione.
    \item Se $A = \mathbb{N}$ non ci sono punti di accumulazione.
    \item Se $A = \mathbb{Q}$ ogni numero reale è di accumulazione per $A$: ogni numero reale si può approssimare con numeri razionali.
\end{itemize}
\end{esempio}

\section{Confronto tra continuità e limite}

Data una funzione $f : D \to \mathbb{R}$ e sia $x_0$ un punto di accumulazione per $D$:
\begin{enumerate}
    \item nella definizione di continuità i valori vengono confrontati con $f(x_0)$; in quella di limite con un valore $\ell$ che può essere diverso da $f(x_0)$;
    \item la condizione di continuità è verificata per $x = x_0$, mentre nella definizione di limite si esclude dal confronto il punto $x = x_0$;
    \item $f$ è continua in $x_0$ equivale a dire che $f$ ha limite $\ell = f(x_0)$ per $x$ tendente a $x_0$;
    \item entrambe le definizioni se sono vere per un certo $\delta > 0$ saranno vere per ogni $\delta' < \delta$.
\end{enumerate}

\section{Esempi}

\begin{esempio}[$f(x) = 1 + x^3$ è continua in $x_0 = 0$]
$D = \mathbb{R}$ e $f(0) = 1 + 0^3 = 1$. Fissato $\varepsilon > 0$:
\[
|f(x) - f(0)| = |1 + x^3 - 1| = |x^3| < \varepsilon \iff |x| < \sqrt[3]{\varepsilon} .
\]
Basta considerare $\delta = \sqrt[3]{\varepsilon}$: per ogni $\varepsilon > 0$, con $\delta = \sqrt[3]{\varepsilon}$, se $|x - x_0| = |x| < \delta = \sqrt[3]{\varepsilon}$ allora $|(1 + x^3) - 1| < \varepsilon$, cioè $|f(x) - f(0)| < \varepsilon$. Dunque $f$ è continua in $x_0 = 0$ e $\lim_{x \to 0} f(x) = f(0)$.
\end{esempio}

\begin{esempio}[Funzione non continua]
Sia
\[
g(x) = \begin{cases} x & x \neq 0 \\ 1 & x = 0 \end{cases}
\qquad D = \mathbb{R} .
\]
\emph{$g$ è continua in $x_0 = 0$?} No. Per ogni $\varepsilon > 0$ dovrebbe esistere $\delta > 0$ tale che per ogni $x \in \mathbb{R}$ con $|x - x_0| = |x| < \delta$ si abbia $|g(x) - g(x_0)| < \varepsilon$. Ma $g(0) = 1$ e per $x \ne 0$ vale $|g(x) - g(x_0)| = |x - 1| < \varepsilon$, cioè $1 - \varepsilon < x < 1 + \varepsilon$. Se scegliamo $\varepsilon$ in modo che $1 - \varepsilon > 0$, per esempio $\varepsilon = \tfrac{1}{3}$, si ha $1 - \varepsilon = \tfrac{2}{3}$ e la condizione diventa $\tfrac{2}{3} < x < \tfrac{4}{3}$. Ma qualsiasi $\delta > 0$ dà $|x| < \delta$, cioè $-\delta < x < \delta$, e nell'intervallo $]-\delta, \delta[$ ci sono dei valori di $x$ per cui $x < \tfrac{2}{3} = 1 - \varepsilon$. Quindi per $\varepsilon = \tfrac{1}{3}$ \emph{non} c'è nessun $\delta > 0$ che soddisfa la definizione di continuità: quando $x$ si avvicina a $0$, $g(x)$ \emph{non} si avvicina a $g(0)$.

\emph{$g$ ha limite per $x \to 0$?} Sì, $\displaystyle\lim_{x\to 0} g(x) = 0$. Verifichiamo con la definizione: per ogni $\varepsilon > 0$ deve esistere $\delta > 0$ tale che per ogni $x \in D \setminus \{x_0\}$ con $|x - x_0| < \delta$ si abbia $|g(x) - \ell| < \varepsilon$. Con $\ell = 0$: $|g(x) - \ell| = |x - 0| < \varepsilon \iff -\varepsilon < x < \varepsilon$. Basta scegliere $\delta = \varepsilon$. Il limite $\ell = 0$ è diverso da $g(0) = 1$.
\end{esempio}

\begin{esempio}[$h(x) = \dfrac{\sin x}{x}$]
Il campo di esistenza è $x \neq 0$, $D = (-\infty, 0) \cup (0, +\infty)$. La funzione $h$ \emph{non è definita} in $x_0 = 0$, quindi non ha senso chiedersi se $h$ è continua in $x_0 = 0$. Tuttavia $x_0 = 0$ è un punto di accumulazione per $D$ e
\[
\lim_{x \to 0} \frac{\sin x}{x} = 1 .
\]
\end{esempio}

\end{document}
